%%% ============================================================
%%%  Experimento: Dataset Aumentado para Predicción S1→S2
%%%  con Búsqueda de Arquitectura NN y Visualización t-SNE
%%%  Compilar con pdfLaTeX.
%%% ============================================================
\documentclass[11pt,a4paper]{article}

\usepackage[utf8]{inputenc}
\usepackage[T1]{fontenc}
\usepackage[spanish]{babel}
\usepackage{amsmath,amssymb}
\usepackage{bm}
\usepackage{booktabs}
\usepackage{graphicx}
\usepackage{hyperref}
\usepackage{geometry}
\usepackage{caption}
\usepackage{subcaption}
\usepackage{multirow}
\usepackage{array}
\usepackage{xcolor}
\usepackage{float}
\usepackage{siunitx}

\geometry{margin=2.5cm}

\sisetup{
  round-mode      = places,
  round-precision = 4,
  table-number-alignment = center,
}

\hypersetup{colorlinks=true, linkcolor=blue, citecolor=blue, urlcolor=blue}

\title{%
    Experimento: Entrenamiento con Dataset Aumentado \\
    y Búsqueda de Arquitectura de Red de Proyección \\
    para Predicción $S_1 \to S_2$
}
\author{Luis Plaza}
\date{\today}

\begin{document}
\maketitle
\tableofcontents
\newpage

% ============================================================
\section{Motivación}
% ============================================================

Los experimentos anteriores de predicción $S_1 \to S_2$ (segunda sección del informe principal) utilizaron una o más cohortes completas de entrenamiento. El mejor resultado documentado —modelo con \textbf{3~cohortes balanceadas, 4D}, $d_r=11$— alcanzó Acc$=0.755$, Recall$=0.548$, Precisión$=0.341$ y F1$\approx 0.421$ sobre la cohorte 2022.

Una limitación de esa estrategia es que la mayoría de los alumnos que aprueban aportan poca señal discriminativa para la clase \texttt{reprueba}. Cuando se mezclan 3 cohortes completas, la clase aprobatoria domina el espacio de embeddings. Este experimento propone una estrategia alternativa de construcción del dataset que conserva \emph{toda} la señal de riesgo histórica sin dilutirla con alumnos que nunca reprueban.

% ============================================================
\section{Dataset Aumentado}
% ============================================================

\subsection{Estrategia de construcción}

El \textbf{dataset aumentado} (\texttt{dataset\_augmentado\_1920reprueba\_2021\_s2}) combina:

\begin{itemize}
    \item \textbf{Cohorte 2021 completa:} todos los alumnos con los 4 cursos de $S_1$
    (BT1211, FI1000, MA1001, MA1101) y sus relaciones de $S_2$
    (CC1002, FI1100, MA1002, MA1102).
    \item \textbf{Cohortes 2019 y 2020 con $\geq 1$ reprobación:} de estas cohortes se
    incorporan \emph{todos} los registros de los alumnos que reprobaron al menos una
    asignatura de $S_2$, incluyendo también sus tripletas de aprobación para conservar
    el contexto relacional.
\end{itemize}

\noindent Esta selección concentra la información histórica de riesgo sin agregar el ruido proveniente de cohortes previas de alumnos puramente aprobatorios.

\subsection{Estadísticas del dataset}

\begin{table}[H]
\centering
\caption{Estadísticas del dataset aumentado versus dataset con 3~cohortes completas.}
\label{tab:stats_dataset}
\begin{tabular}{lcc}
\toprule
\textbf{Estadístico} & \textbf{3 Cohortes completas} & \textbf{Dataset aumentado} \\
\midrule
Total de tripletas & ${\sim}$8\,000 & 3\,606 \\
Entidades (alumnos + cursos) & ${\sim}$2\,000 & 959 \\
Alumnos en entrenamiento NN & ${\sim}$2\,000 & 951 \\
Cohortes de entrenamiento & 2019, 2020, 2021 & 2021 + rep.\ 2019/2020 \\
\bottomrule
\end{tabular}
\end{table}

\subsection{Protocolo de entrenamiento TuckER}

Se entrenaron modelos TuckER para $d_r \in \{1, \ldots, 12\}$ con los siguientes hiperparámetros:

\begin{itemize}
    \item Dimensión de embeddings de entidad: $d_e = 4$.
    \item Inicialización \emph{warm-start}: cada embedding de alumno se inicializa con
    $(n_c - 4.0)/3.0$ para $c \in \{$BT1211, FI1000, MA1001, MA1101$\}$, llevando el
    rango $[1, 7]$ al intervalo $[-1, 1]$.
    \item Optimizador: Adam, lr $= 0.003$.
    \item Dropout: $(0.3,\,0.4,\,0.5)$ en entrada y capas ocultas.
    \item Épocas máximas: 500, con \emph{early stopping} por MRR de validación
    (paciencia 100).
\end{itemize}

Las cohortes OOT son exclusivamente $20221 \to 20222$ y $20231 \to 20232$ (las cohortes
2019, 2020 y 2021 están contenidas en el entrenamiento).

% ============================================================
\section{Búsqueda de Arquitectura de Red de Proyección}
% ============================================================

\subsection{Metodología}

La red de proyección $f_\theta:\mathbb{R}^4 \to \mathbb{R}^{d_e}$ mapea el vector de notas
normalizadas de $S_1$ hacia el embedding aprendido por TuckER.  Se evaluaron 23
configuraciones que varían:

\begin{itemize}
    \item \textbf{Topología de capas ocultas:} desde una sola capa de 16 neuronas hasta
    arquitecturas profundas simétricas como $[64, 128, 64]$.
    \item \textbf{Dropout:} $\{0.0, 0.1, 0.2, 0.3, 0.5\}$.
    \item \textbf{Función de activación:} ReLU, Tanh, LeakyReLU($\alpha=0.1$).
    \item \textbf{Tasa de aprendizaje:} $\{10^{-3},\,5 \times 10^{-4},\,10^{-4}\}$.
\end{itemize}

Cada configuración se entrenó con SEED$=42$, máximo 500 épocas y \emph{early stopping}
(paciencia 100) con pérdida MSE. Los modelos OOT se evaluaron sobre
$20221 \to 20222$ y $20231 \to 20232$, reportando F1 y Recall de la clase
\texttt{reprueba}.

\subsection{Resultados para $d_r = 5$}

\begin{table}[H]
\centering
\caption{TOP-5 configuraciones por F1 promedio (OOT) para $d_r=5$.}
\label{tab:top5_f1_rdim5}
\begin{tabular}{lcccccc}
\toprule
\textbf{Config} & \textbf{Arquitectura} & \textbf{Drop} & \textbf{Act} & \textbf{F1$_{avg}$} & \textbf{Rec$_{avg}$} & \textbf{Épocas} \\
\midrule
\texttt{drop05}      & [64, 128]    & 0.5 & ReLU   & \textbf{0.3985} & 0.7207 & 300 \\
\texttt{leaky\_wider} & [128, 256]  & 0.2 & Leaky  & 0.3980 & 0.7502 & 105 \\
\texttt{tanh\_deep}  & [64, 128, 64] & 0.2 & Tanh  & 0.3967 & 0.7115 & 479 \\
\texttt{leaky\_base} & [64, 128]   & 0.3 & Leaky  & 0.3952 & 0.7575 & 161 \\
\texttt{drop0}       & [64, 128]   & 0.0 & ReLU   & 0.3913 & 0.7630 & 111 \\
\midrule
\textit{baseline}    & [64, 128]   & 0.3 & ReLU   & 0.3642 & 0.8245 & 127 \\
\bottomrule
\end{tabular}
\end{table}

\begin{table}[H]
\centering
\caption{TOP-5 configuraciones por Recall promedio (OOT) para $d_r=5$.}
\label{tab:top5_rec_rdim5}
\begin{tabular}{lcccccc}
\toprule
\textbf{Config} & \textbf{Arquitectura} & \textbf{Drop} & \textbf{Act} & \textbf{F1$_{avg}$} & \textbf{Rec$_{avg}$} & \textbf{Épocas} \\
\midrule
\texttt{large}      & [256, 512]   & 0.4 & ReLU  & 0.3829 & \textbf{0.8489} & 104 \\
\texttt{deep\_wide} & [128, 256, 128] & 0.3 & ReLU & 0.3407 & 0.8401 & 132 \\
\texttt{lr1e4}      & [64, 128]   & 0.3 & ReLU  & 0.3638 & 0.8352 & 289 \\
\texttt{baseline}   & [64, 128]   & 0.3 & ReLU  & 0.3642 & 0.8245 & 127 \\
\texttt{wider}      & [128, 256]  & 0.3 & ReLU  & 0.3635 & 0.8189 & 111 \\
\bottomrule
\end{tabular}
\end{table}

\begin{table}[H]
\centering
\caption{Desglose por cohorte OOT para la mejor configuración por F1 y por Recall ($d_r=5$).}
\label{tab:detalle_rdim5}
\begin{tabular}{llcccc}
\toprule
\textbf{Config (objetivo)} & \textbf{Cohorte} & \textbf{F1} & \textbf{Prec} & \textbf{Rec} & \textbf{Acc} \\
\midrule
\multirow{2}{*}{\texttt{drop05} (F1)}
  & $20221 \to 20222$ & 0.3783 & 0.2669 & 0.6488 & 0.6543 \\
  & $20231 \to 20232$ & 0.4188 & 0.2846 & 0.7926 & 0.6317 \\
\midrule
\multirow{2}{*}{\texttt{large} (Recall)}
  & $20221 \to 20222$ & 0.3802 & 0.2494 & 0.7996 & 0.5775 \\
  & $20231 \to 20232$ & 0.3855 & 0.2454 & 0.8981 & 0.5208 \\
\bottomrule
\end{tabular}
\end{table}

\subsection{Resultados para $d_r = 8$}

\begin{table}[H]
\centering
\caption{TOP-5 configuraciones por F1 promedio (OOT) para $d_r=8$.}
\label{tab:top5_f1_rdim8}
\begin{tabular}{lcccccc}
\toprule
\textbf{Config} & \textbf{Arquitectura} & \textbf{Drop} & \textbf{Act} & \textbf{F1$_{avg}$} & \textbf{Rec$_{avg}$} & \textbf{Épocas} \\
\midrule
\texttt{deep\_small} & [32, 64, 32]  & 0.2 & ReLU  & \textbf{0.4023} & 0.5546 & 295 \\
\texttt{deep\_sym}   & [64, 128, 64] & 0.3 & ReLU  & 0.3814 & 0.5824 & 403 \\
\texttt{combo\_tanh\_lr} & [128, 256] & 0.2 & Tanh & 0.3385 & 0.6104 & 357 \\
\texttt{tanh\_deep}  & [64, 128, 64] & 0.2 & Tanh  & 0.3345 & 0.6488 & 284 \\
\texttt{drop05}      & [64, 128]    & 0.5 & ReLU  & 0.3339 & 0.5963 & 393 \\
\midrule
\textit{baseline}    & [64, 128]    & 0.3 & ReLU  & 0.1929 & 0.2076 & 119 \\
\bottomrule
\end{tabular}
\end{table}

\begin{table}[H]
\centering
\caption{TOP-5 configuraciones por Recall promedio (OOT) para $d_r=8$.}
\label{tab:top5_rec_rdim8}
\begin{tabular}{lcccccc}
\toprule
\textbf{Config} & \textbf{Arquitectura} & \textbf{Drop} & \textbf{Act} & \textbf{F1$_{avg}$} & \textbf{Rec$_{avg}$} & \textbf{Épocas} \\
\midrule
\texttt{tanh\_deep}      & [64, 128, 64] & 0.2 & Tanh & 0.3345 & \textbf{0.6488} & 284 \\
\texttt{combo\_tanh\_lr} & [128, 256]    & 0.2 & Tanh & 0.3385 & 0.6104 & 357 \\
\texttt{tiny}            & [16, 32]     & 0.1 & ReLU & 0.3275 & 0.6073 & 329 \\
\texttt{single\_128}     & [128]        & 0.2 & ReLU & 0.2841 & 0.6063 & 145 \\
\texttt{drop05}          & [64, 128]    & 0.5 & ReLU & 0.3339 & 0.5963 & 393 \\
\bottomrule
\end{tabular}
\end{table}

\begin{table}[H]
\centering
\caption{Desglose por cohorte OOT para la mejor configuración por F1 y por Recall ($d_r=8$).}
\label{tab:detalle_rdim8}
\begin{tabular}{llcccc}
\toprule
\textbf{Config (objetivo)} & \textbf{Cohorte} & \textbf{F1} & \textbf{Prec} & \textbf{Rec} & \textbf{Acc} \\
\midrule
\multirow{2}{*}{\texttt{deep\_small} (F1)}
  & $20221 \to 20222$ & 0.3941 & 0.3053 & 0.5556 & 0.7232 \\
  & $20231 \to 20232$ & 0.4104 & 0.3261 & 0.5537 & 0.7337 \\
\midrule
\multirow{2}{*}{\texttt{tanh\_deep} (Recall)}
  & $20221 \to 20222$ & 0.3369 & 0.2263 & 0.6587 & 0.5797 \\
  & $20231 \to 20232$ & 0.3321 & 0.2243 & 0.6389 & 0.5697 \\
\bottomrule
\end{tabular}
\end{table}

\subsection{Observaciones sobre la búsqueda de arquitectura}

El comportamiento de $d_r$ frente a la arquitectura NN es marcadamente diferente:

\begin{itemize}
    \item \textbf{$d_r=5$:} favorece regularización fuerte.
    \texttt{drop05} (dropout$=0.5$) y configuraciones \texttt{leaky} dominan en F1.
    Las arquitecturas grandes (\texttt{large}, \texttt{deep\_wide}) dominan en Recall
    a costa de precisión. Las architecturas de baja lr ($10^{-4}$) también son
    competitivas en Recall, sugiriendo que el espacio de dimensión 5 requiere un
    proyector que converja lentamente para no sobreajustar los embeddings de entrenamiento.
    \item \textbf{$d_r=8$:} el espacio de relaciones más rico requiere un proyector
    \emph{compacto y profundo}. \texttt{deep\_small}$=[32, 64, 32]$ es el claro ganador en F1;
    las arquitecturas grandes (\texttt{baseline} $[64, 128]$, \texttt{large}, \texttt{wider})
    colapsan casi completamente (F1$\approx 0.19$). La arquitectura base por defecto
    resulta inútil para $d_r=8$, mientras que para $d_r=5$ es moderadamente competitiva.
\end{itemize}

% ============================================================
\section{Evaluación OOT: Mejora respecto a arquitectura por defecto}
% ============================================================

La tabla siguiente compara la arquitectura por defecto (baseline, $[64,128]$, drop$=0.3$)
con la mejor arquitectura encontrada para cada $d_r$.

\begin{table}[H]
\centering
\caption{Comparación entre arquitectura por defecto y mejor arquitectura NN
para el dataset aumentado, sobre las dos cohortes OOT.
Métricas referidas a la clase \texttt{reprueba}.}
\label{tab:baseline_vs_best}
\begin{tabular}{llcccc}
\toprule
\textbf{Modelo} & \textbf{Cohorte OOT} & \textbf{F1} & \textbf{Prec} & \textbf{Rec} & \textbf{Acc} \\
\midrule
\multicolumn{6}{l}{\textit{$d_r=5$ — arquitectura por defecto [64,128], drop=0.3}} \\
Baseline ($d_r=5$) & $20221 \to 20222$ & 0.3595 & 0.2327 & 0.7897 & 0.5441 \\
Baseline ($d_r=5$) & $20231 \to 20232$ & 0.3439 & 0.2129 & 0.8944 & 0.4287 \\
\midrule
\multicolumn{6}{l}{\textit{$d_r=5$ — mejor config: \texttt{drop05} [64,128], drop=0.5}} \\
\texttt{drop05} ($d_r=5$) & $20221 \to 20222$ & 0.3783 & 0.2669 & 0.6488 & 0.6543 \\
\texttt{drop05} ($d_r=5$) & $20231 \to 20232$ & \textbf{0.4188} & 0.2846 & 0.7926 & 0.6317 \\
\midrule
\multicolumn{6}{l}{\textit{$d_r=8$ — arquitectura por defecto [64,128], drop=0.3}} \\
Baseline ($d_r=8$) & $20221 \to 20222$ & 0.3429 & 0.2692 & 0.4722 & 0.7068 \\
Baseline ($d_r=8$) & $20231 \to 20232$ & 0.4062 & 0.3273 & 0.5352 & 0.7381 \\
\midrule
\multicolumn{6}{l}{\textit{$d_r=8$ — mejor config: \texttt{deep\_small} [32,64,32], drop=0.2}} \\
\texttt{deep\_small} ($d_r=8$) & $20221 \to 20222$ & 0.3941 & 0.3053 & 0.5556 & 0.7232 \\
\texttt{deep\_small} ($d_r=8$) & $20231 \to 20232$ & \textbf{0.4104} & 0.3261 & 0.5537 & 0.7337 \\
\bottomrule
\end{tabular}
\end{table}

\paragraph{Síntesis.}
Para $d_r=5$, la búsqueda de arquitectura mejora el F1 promedio en
$+0.047$ ($+13\%$), elevándolo de $0.352$ a $0.399$. La precisión sube notoriamente
($0.233 \to 0.276$) mientras el Recall baja moderadamente ($0.842 \to 0.721$): las
predicciones resultan más fiables aunque menos exhaustivas.
Para $d_r=8$, la mejora en F1 es $+0.028$ ($+7\%$) y el Recall también sube ligeramente
($0.504 \to 0.555$), indicando que en este caso la arquitectura correcta ayuda tanto
en precisión como en cobertura.

% ============================================================
\section{Comparación con Modelos del Informe Principal}
% ============================================================

La Tabla~\ref{tab:comparacion_modelos} contrasta los mejores modelos del informe principal
(entrenados con 1 y 3 cohortes completas) con los mejores modelos del dataset aumentado.
Todos los valores corresponden a la cohorte $20221 \to 20222$, que es la cohorte de prueba
en común entre ambos diseños experimentales.

\begin{table}[H]
\centering
\caption{Comparación de modelos de predicción $S_1 \to S_2$ sobre cohorte 2022.
Métricas de la clase \texttt{reprueba}. Los valores del informe provienen de la
Tabla~2 de validación de cohortes (mejor $d_r$ por estrategia). Los valores del
dataset aumentado corresponden a la cohorte OOT $20221\to 20222$.
$^\dagger$ F1 calculado a partir de Prec y Rec reportados.}
\label{tab:comparacion_modelos}
\resizebox{\textwidth}{!}{%
\begin{tabular}{lcccccc}
\toprule
\textbf{Modelo} & \textbf{Train} & \textbf{Acc} & \textbf{Prec} & \textbf{Rec} & \textbf{F1} & \textbf{$d_r$} \\
\midrule
\multicolumn{7}{l}{\textit{Informe principal (evaluados en cohorte 2022)}} \\
1 Cohorte 4D, mejor Prec & 2019 & 0.945 & \textbf{0.598} & 0.315 & 0.413$^\dagger$ & 6 \\
1 Cohorte 5D Bal., mejor F1 & 2019 & 0.736 & 0.117 & 0.511 & 0.190$^\dagger$ & 10 \\
3 Cohortes 4D Bal., mejor & 2019+2020+2021 & 0.755 & 0.341 & 0.548 & 0.421$^\dagger$ & 11 \\
\midrule
\multicolumn{7}{l}{\textit{Dataset aumentado (OOT $20221\to 20222$)}} \\
Aug.\ baseline $d_r=5$         & Aug. & 0.544 & 0.233 & 0.790 & 0.360 & 5 \\
Aug.\ \texttt{drop05} $d_r=5$  & Aug. & 0.654 & 0.267 & 0.649 & 0.378 & 5 \\
Aug.\ baseline $d_r=8$         & Aug. & 0.707 & 0.269 & 0.472 & 0.343 & 8 \\
Aug.\ \textbf{\texttt{deep\_small}} $d_r=8$ & Aug. & 0.723 & 0.305 & 0.556 & \textbf{0.394} & 8 \\
\midrule
\multicolumn{7}{l}{\textit{Dataset aumentado (OOT $20231\to 20232$, sin equivalente en informe)}} \\
Aug.\ \texttt{drop05} $d_r=5$      & Aug. & 0.632 & 0.285 & 0.793 & 0.419 & 5 \\
Aug.\ \textbf{\texttt{deep\_small}} $d_r=8$ & Aug. & 0.734 & 0.326 & 0.554 & \textbf{0.410} & 8 \\
\bottomrule
\end{tabular}%
}
\end{table}

\paragraph{Lectura de la comparación.}
Tres observaciones principales emergen de esta comparación:

\begin{enumerate}
    \item \textbf{El modelo de 1~cohorte (4D, $d_r=6$)} tiene la mayor precisión ($0.598$),
    pero un Recall muy bajo ($0.315$): detecta pocos reprobados pero con alta confianza.
    No se generaliza bien a otras cohortes y su F1 ($0.413$) es comparable al del modelo
    aumentado \texttt{deep\_small} ($0.394$).
    
    \item \textbf{El mejor modelo del informe} (3~cohortes balanceadas, 4D, $d_r=11$)
    logra F1$=0.421$ con Recall$=0.548$ y Prec$=0.341$. El modelo aumentado
    \texttt{deep\_small} ($d_r=8$) obtiene F1$=0.394$ sobre la misma cohorte 2022,
    con Recall similar ($0.556$) pero con $\approx 13\%$ más de datos de entrenamiento
    \emph{menos}: el dataset aumentado usa solo alumnos de riesgo de 2019/2020 más la
    cohorte 2021 completa, versus 3~cohortes completas.
    
    \item \textbf{La ventaja del dataset aumentado} se manifiesta en su capacidad de
    generalización a una cohorte adicional ($20231\to 20232$), donde el modelo
    \texttt{drop05} $d_r=5$ alcanza F1$=0.419$, superando el mejor resultado del informe
    sobre la misma tarea. Esto sugiere que la estrategia de selección de datos de riesgo
    histórico es más eficiente y estable que agregar cohortes completas.
\end{enumerate}

\paragraph{Nota metodológica.}
Los modelos del informe principal fueron evaluados sobre todos los cursos de $S_2$ presentes
en la cohorte 2022; los modelos del dataset aumentado se evalúan sobre los cuatro cursos
fundamentales de $S_2$: CC1002, FI1100, MA1002, MA1102. Esta diferencia de alcance debe
tenerse en cuenta al interpretar la comparación.

% ============================================================
\section{Visualización t-SNE del Espacio Latente}
% ============================================================

Para comprender cómo el modelo aprende a separar alumnos en riesgo de los que no lo están,
se proyectó el vector de \emph{perfil relacional} $v_{s,o} \in \mathbb{R}^{d_r}$ de cada
alumno OOT hacia 2D mediante t-SNE. Este vector se define como:
\[
v_{s,o} = \mathcal{W} \times_2 \hat{e}_s \times_3 e_o
\]
donde $\mathcal{W} \in \mathbb{R}^{d_r \times d_e \times d_e}$ es el tensor núcleo de TuckER,
$\hat{e}_s = f_\theta(x_s)$ es el embedding del alumno proyectado por la red neuronal y
$e_o$ es el embedding del curso objetivo aprendido por TuckER.

La proyección se realizó para el modelo $d_r=8$ con arquitectura \texttt{deep\_small}
sobre la cohorte OOT $20231 \to 20232$.

\subsection{Separación aprueba/reprueba}

\begin{figure}[H]
    \centering
    \includegraphics[width=\textwidth]{../results/tsne_augmentado_rdim8/tsne_binario_rdim8_2023120232.png}
    \caption{t-SNE de perfiles relacionales para los 4 cursos de $S_2$, cohorte
    OOT $20231 \to 20232$, modelo aumentado $d_r=8$ (\texttt{deep\_small}).
    Los puntos en \textcolor{red}{rojo} corresponden a alumnos que reprobaron, los
    \textcolor{green}{verdes} a quienes aprobaron.}
    \label{fig:tsne_binario}
\end{figure}

En CC1002 y MA1002 se observan agrupaciones diferenciadas de reprobados (esquina derecha
en CC1002, zona izquierda en MA1002), lo que indica que el modelo aprende regiones latentes
con mayor concentración de riesgo. FI1100 y MA1102 muestran mayor mezcla, consistente con
el bajo número de reprobaciones ($n=26$ y $n=18$ respectivamente) que dificultan
la formación de fronteras discriminativas.

\subsection{Gradiente de nota final}

\begin{figure}[H]
    \centering
    \includegraphics[width=\textwidth]{../results/tsne_augmentado_rdim8/tsne_notas_rdim8_2023120232.png}
    \caption{t-SNE de perfiles relacionales coloreado por nota final (escala continua
    rojo-amarillo-verde, $1.0 \to 7.0$), mismas condiciones que la Figura~\ref{fig:tsne_binario}.}
    \label{fig:tsne_notas}
\end{figure}

La coloración continua por nota revela que el modelo aprende un gradiente de rendimiento
\emph{suave}: en MA1002 se aprecia un gradiente espacial de notas bajas (izquierda) a
notas altas (derecha), lo que indica que el espacio latente codifica no solo la
distinción binaria sino la magnitud del rendimiento. Este comportamiento es especialmente
notorio en MA1002 y, en menor medida, en CC1002.

% ============================================================
\section{Conclusiones}
% ============================================================

\begin{enumerate}
    \item El dataset aumentado (cohorte 2021 completa + reprobados 2019/2020) es una
    estrategia eficiente de construcción de datos: con ${\sim}60\%$ menos tripletas que
    el dataset de 3~cohortes completas, alcanza métricas OOT comparables o superiores.
    
    \item La búsqueda de arquitectura de la red de proyección es \textbf{crítica}:
    la arquitectura por defecto $[64,128]$ es inútil para $d_r=8$ (F1$=0.193$),
    mientras que \texttt{deep\_small} $[32,64,32]$ eleva el F1 a $0.402$.
    Para $d_r=5$, el beneficio es menor pero consistente ($+13\%$ en F1).
    
    \item El tradeoff entre precisión y recall depende del $d_r$:
    $d_r=8$ + \texttt{deep\_small} ofrece mayor precisión (Prec$\approx 0.31$, menor ruido),
    mientras que $d_r=5$ + \texttt{drop05} ofrece mayor recall (Rec$\approx 0.72$, menor
    omisión). La elección depende del contexto de uso del sistema de alerta.
    
    \item Las visualizaciones t-SNE confirman que el espacio latente captura estructura
    semántica real: los alumnos en riesgo se agrupan en regiones diferenciadas para los
    cursos con mayor base de reprobación (CC1002, MA1002), y existe un gradiente
    continuo de rendimiento alineado con la dimensión espacial.
    
    \item La comparación con el informe principal muestra que el modelo aumentado supera
    el mejor resultado histórico (3~cohortes, F1$=0.421$) en la cohorte $20231 \to 20232$
    con \texttt{drop05} $d_r=5$ (F1$=0.419$), siendo además evaluable en una cohorte
    adicional gracias a que el diseño OOT excluye las cohortes de entrenamiento.
\end{enumerate}

\end{document}
